\section{数据与方法}
\subsection{研究窗口与历史样本}
三类股票池共同使用2007-01至2025-12信号月，共228期。实际建仓从2007-02-01收盘开始，最后退出为2026-02-02收盘；持有期按建仓月份归属2007-02至2026-01。2026-01-30仅用于最新完整截面的分布诊断，其计划退出日超出行情截止，未进入绩效。

全部 A 股沿用课程沪深京人民币权益证券覆盖，包含库内境内上市中国存托凭证（CDR）；北交所自2021-11-15开市后纳入。历史生命周期、指数成员、行业和ST/PT状态均按有效日期匹配，不用当前名单或市值排名代替历史指数池。入场资格采用下一交易日收盘可交易口径，全天或收盘仍停牌者剔除，日内临停后已恢复收盘交易者可保留。

三个池分别在S上进行预处理和分层，之后才接入未来响应Y。完整窗口的全部 A 股共有686,280个股票--期。原始8,011个退出端点不足中，7,936项按停牌与公司行为规则估值，58项依据事先公告在最后实际交易日退出后持有零利息现金，17项按合并对价估值。其中两项承接A股尚未上市，采用已确认的换股权利估值延续。处理后未解决端点为0，原持仓与计划持有期未因未来可观测性改变。

行业使用历史生效事件，保留早期分类，不把现行分类倒填。229个诊断月合计87个F股票--期缺行业，按预先确定的控制变量规则从S排除。有限流通市值异常已按公司股本事件修正。这些修复支持当前历史字段研究，不构成全市场财报、行业与股本版本的严格PIT重建。

\subsection{默认测试设定}
普通最小二乘法（OLS）用于因子中性化，加权最小二乘法（WLS）用于收益回归；无风险收益率（RF）设为0。

\begin{table}[H]\centering
\caption*{\textbf{表2　统一参数与评价定义}}
\small\renewcommand{\arraystretch}{1.18}
\begin{tabularx}{\linewidth}{@{}p{25mm}X@{}}\toprule
项目 & 本报告采用的设定 \\ \midrule
频率与时点 & 月末生成信号；下一交易日收盘换仓，至下一次换仓日收盘结束 \\
去极值 & 截面中位数$\pm5\times\mathrm{MAD}$，MAD为绝对离差的中位数 \\
回归暴露 & 去极值后的标准分数（Z-score），截面标准差分母为$n$（ddof=0） \\
中性化 & 去极值BP对截距、$\ln(\text{总市值})$及行业哑变量作OLS，取残差 \\
WLS & 下期收益对标准化BP、对数总市值和行业回归；目标函数权重$w_i=\sqrt{\text{流通市值}_i}$ \\
Rank IC & 冻结的中性化残差与下期个股收益的Spearman相关 \\
分层 & 五层等人数、层内等权；并列按证券代码稳定排序；十层与分域复用相同残差 \\
成本 & 教材只数换手，首期为1；每层扣$0.004\times TO$；多空两腿分别扣费 \\
内部基准 & 同期全部实际分层股票的等权毛收益，不另扣换手成本 \\
外部指数 & 沪深300、中证500价格指数；共同入/退出端点，不含分红 \\
时间统计 & RF=0；时间序列标准差分母为$n-1$；ICIR不年化 \\
缺值与缺期 & 不填补因子；无法合理估值的持仓不重分权；缺月不补零或拼接年化 \\
\bottomrule\end{tabularx}
\end{table}

回归检验线性条件关联，Rank IC检验截面排序，分层则检验离散组合的成本后收益。三者使用相同日期与形成样本，但权重和关系形状不同，不要求数值等价。负的中性化残差仍是有效因子值。退化截面、设计矩阵不满秩或样本不足均显式记录；本次三池主窗口均可完整评价。

个股收益由同一原始代码复权尺度的入、退出估值构造。教材只数换手与多空净收益分别为
\begin{align}
TO_{q,t}&=1-\frac{|S_{q,t}\cap S_{q,t-1}|}{|S_{q,t}|},\qquad TO_{q,1}=1,\\
r^{LS,net}_t&=r^{gross}_{5,t}-r^{gross}_{1,t}-0.004TO_{5,t}-0.004TO_{1,t}.
\end{align}
这一近似不模拟权重漂移、完整再平衡、成交冲击或融资借券。多空净值按$\prod_t(1+r^{LS,net}_t)$计算，沿用教材资本口径，不另除以2。

对连续完整月度净收益$r_t$，复合年化收益率（CAGR）为$[\prod_t(1+r_t)]^{12/n}-1$，夏普比率为$\sqrt{12}\,\bar r/s_r$。设月超额$e_t=r_{q,t}^{net}-r_{b,t}^{gross}$，算术年化超额为$12\bar e$，信息比率（IR）为$\sqrt{12}\,\bar e/s_e$。净值均含期初1；最大回撤为净值相对历史峰值的最小变化，报告为非正数。超额最大回撤对相对净值$NAV_q/NAV_b$作相同计算。

C1的信息系数比率（ICIR）为月IC均值除以样本标准差，不乘$\sqrt{12}$。截面$t$使用课程常规WLS方差估计，未作稳健标准误调整；其均值和绝对值均值不构成因子收益时间序列均值的显著性检验。D2与D3分别累加IC和回归系数，不复利。D6在每个连续12期窗口内计算层收益复利与基准收益复利之差，前11期为空。

\subsection{分域比较与解释窗口}
D9在全部 A 股同一形成样本与残差上重分十层。D10每月先按流通市值等人数分为低、中、高三域，再在域内按同一残差十等分。C4在历史一级行业内部重分五层，不重新估计行业内规模斜率。分域基准均为域内实际分层股票的等权毛收益；控制变量缺失者不为构造基准重新纳入。这是相对附录“因子非空”定义的明确资格限制。

C4每期至少25只，D10每域至少50只，并检验不同因子取值数。C4共同信号窗在查看收益前冻结为2021-07至2025-12，共54期，实际入场2021-08-02、退出2026-02-02。综合行业只有12期达到门槛，保留不足行，不拼接年化或进入同窗排名。

时间解读采用三个连续日历块：2007-02至2013-12、2014-01至2020-12、2021-01至2026-01，均按持有期起始月归属。这些分段仅作描述，不用于调整主窗口或额外推断显著性。内部基准与分层的成本待遇不同，外部价格指数与个股复权收益的分红口径也不同，后文分别解释。
